\section{Related Work}

\paragraph{Learning-rate stability and the Hessian paradigm.}
Classical convergence and stability analyses of gradient descent are typically
stated in terms of smoothness constants or Hessian curvature, leading to
step-size constraints proportional to $1/\lambda_{\max}(H)$.
This Hessian-centric viewpoint underlies much of modern optimization theory
for machine learning and motivates the common practice of diagnosing training
instabilities through curvature (sharpness) measurements.

\paragraph{Edge of Stability (EoS) and oscillatory optimization dynamics.}
Cohen et al.~\cite{cohen2021eos} empirically identified that full-batch gradient descent
in neural network training often operates in an \emph{Edge-of-Stability} regime,
where the top Hessian eigenvalue fluctuates near $2/\eta$ and the loss exhibits
short-term oscillations while decreasing over long horizons.
Arora et al.~\cite{arora2022eos} provided a theoretical analysis of mechanisms by which
gradient-based methods can enter and make progress in such EoS regimes,
thereby initiating a line of work on understanding learning dynamics beyond
the classical monotone-descent picture.
More recently, Cohen et al.~\cite{cohen2024centralflows} proposed \emph{central flows},
a differential-equation formalism aimed at modeling the time-averaged trajectory
of oscillatory optimizers in EoS-like regimes, providing a complementary tool
for reasoning about long-term dynamics.

\paragraph{Beyond GD: EoS-style analyses for alternative update rules.}
Extensions of the EoS perspective have also appeared for algorithms that modify
the objective or update rule.
For example, Long and Bartlett~\cite{long2024sameos} derive an ``edge of stability''
condition for Sharpness-Aware Minimization (SAM) and empirically validate that
SAM operates near this boundary; see also Foret et al.~\cite{foret2020sam} for
the original SAM formulation.

\paragraph{Matrix-aware preconditioning and matrix-whitening optimizers.}
A separate but related thread studies \emph{structure-aware} optimizers that exploit
the matrix structure of gradients rather than treating parameters as vectors.
PolarGrad~\cite{lau2025polargrad} provides a unifying perspective for ``matrix-aware''
preconditioning and frames Muon as a special instance (via polar decomposition / related
matrix transforms), with both theoretical and empirical evidence of acceleration.
In parallel, practical scaling work has shown that Muon can be made viable for large-scale
language-model training through techniques such as weight decay and per-parameter scaling~\cite{liu2025muonscalable}.
At the same time, the rapidly growing space of matrix-whitening methods has motivated
careful ablations; Frans et al.~\cite{frans2025matrixwhitening} argue that performance
differences among whitening-style optimizers are not explained solely by the accuracy
of spectral normalization, highlighting the role of variance adaptation.

\paragraph{Spike--tail spectral geometry and spike-aware spectral shaping.}
Motivated by spectral anisotropy in large language model training, Huang et al.~\cite{huang2026spectra}
propose Spectra, a spike-aware optimizer that tracks and suppresses a dominant low-rank spike
subspace without amplifying the spectral tail.
Our work is complementary in spirit to these optimizer developments but is primarily
\emph{learning-theoretic}: we focus on deriving stability and convergence properties of
gradient dynamics under a domain-motivated spike--tail Jacobian structure (including mean bias),
and then interpret normalization, mean-removal, and whitening-style interventions
as geometric operations that reduce dominant singular components and relax step-size constraints.